\documentclass[10pt,a4paper,twocolumn]{article}

% ─── PAQUETES ───────────────────────────────────────────────
\usepackage[utf8]{inputenc}
\usepackage[spanish]{babel}
\usepackage[T1]{fontenc}
\usepackage{geometry}
\geometry{margin=1.9cm, top=1.9cm, bottom=1.9cm, columnsep=0.7cm}
\usepackage{graphicx}
\usepackage{float}
\usepackage{booktabs}
\usepackage{amsmath}
\usepackage{hyperref}
\usepackage{xcolor}
\usepackage{listings}
\usepackage{caption}
\usepackage{subcaption}
\usepackage{fancyhdr}
\usepackage{titlesec}
\usepackage{setspace}
\usepackage{array}
\usepackage{multirow}
\usepackage{url}
\usepackage{enumitem}
\usepackage{parskip}
\usepackage{lastpage}
\usepackage{abstract}
\usepackage{titling}
\usepackage{colortbl}

% ─── ESTILO ─────────────────────────────────────────────────
\setlength{\parindent}{0pt}
\setlength{\parskip}{0.22cm plus 0.08cm minus 0.06cm}
\renewcommand{\baselinestretch}{1.02}

\pagestyle{fancy}
\fancyhf{}
\fancyhead[L]{\small GeoVision-CLIP Cali · Informe Final}
\fancyhead[R]{\small Analítica de Datos I · UAO}
\fancyfoot[C]{\thepage}
\renewcommand{\headrulewidth}{0.3pt}
\setlength{\headheight}{13pt}

% Secciones más aireadas
\titleformat{\section}[hang]{\normalsize\bfseries}{\thesection}{0.5em}{}
\titlespacing{\section}{0pt}{0.4cm plus 0.15cm}{0.2cm plus 0.08cm}
\titleformat{\subsection}[hang]{\small\bfseries}{\thesubsection}{0.5em}{}
\titlespacing{\subsection}{0pt}{0.28cm plus 0.08cm}{0.12cm plus 0.04cm}
\titleformat{\subsubsection}[hang]{\small\bfseries\itshape}{\thesubsubsection}{0.5em}{}
\titlespacing{\subsubsection}{0pt}{0.2cm plus 0.05cm}{0.08cm plus 0.02cm}

\captionsetup{font={small,stretch=0.98},labelfont=bf,skip=5pt}

\definecolor{codeblue}{rgb}{0.0, 0.3, 0.6}
\definecolor{codegreen}{rgb}{0.0, 0.5, 0.0}
\lstset{
  basicstyle=\scriptsize\ttfamily,
  keywordstyle=\color{codeblue}\bfseries,
  commentstyle=\color{codegreen},
  breaklines=true,
  frame=single,
  numbers=left,
  numberstyle=\tiny,
  language=Python,
  tabsize=2,
  showstringspaces=false,
  xleftmargin=0.3cm,
  xrightmargin=0.3cm
}

% Colores para tablas de confiabilidad
\definecolor{greenconf}{rgb}{0.75, 0.95, 0.75}
\definecolor{yellowconf}{rgb}{1.0, 1.0, 0.75}
\definecolor{orangeconf}{rgb}{1.0, 0.85, 0.65}
\definecolor{redconf}{rgb}{1.0, 0.75, 0.75}

\hypersetup{colorlinks=true,linkcolor=blue!60!black,urlcolor=blue!60!black,citecolor=blue!60!black}

% ─── TÍTULO Y AUTORES ──────────────────────────────────────
\title{\vspace{-0.7cm}\Large\bfseries GeoVision-CLIP Cali\\[0.15cm]
       \large Estimación de Contaminación Atmosférica en Puntos No Muestreados\\[0.05cm]
       \normalsize mediante Deep Learning + Estadística Geoespacial Avanzada}

\author{\small
  Nicolás Peña Irurita$^{1}$\hspace{0.3cm}Samuel Patiño Lucumí$^{1}$\\[0.1cm]
  \footnotesize $^{1}$Universidad Autónoma de Occidente, Facultad de Ingenierías, Cali, Colombia\\[0.1cm]
  \footnotesize Programa de Ingeniería de Datos e IA · Tercer Corte · Unidad III · Analítica de Datos I\\[0.05cm]
  \footnotesize Docentes: Prof. Carlos Ferro · Prof. Cristian E. García\\[0.05cm]
  \footnotesize \url{https://github.com/nicothinn/GeoVision-CLIP-Cali}
}
\date{\footnotesize \today}

% ─── INICIO DEL DOCUMENTO ──────────────────────────────────
\begin{document}

    \maketitle
    \begin{abstract}
    Este documento presenta la construcción de un panel analítico espacio-temporal y la arquitectura cloud para la estimación de contaminación atmosférica en el área metropolitana de Santiago de Cali, Colombia. El panel integra cinco años (2019--2023) de datos satelitales multi-fuente---incluyendo gases traza de Sentinel-5P TROPOMI, covariables ópticas de Sentinel-2 MSI, reanálisis meteorológico ERA5, profundidad óptica de aerosoles MODIS MCD19A2---y cuatro años (2020--2023) de mediciones \textit{in-situ} de la red SVCASC (Sistema de Vigilancia de Calidad del Aire de Santiago de Cali). Los datos DAGMA/SVCASC cubren únicamente el período 2020--2023 debido a que el sistema de monitoreo generó predicciones retrospectivas a partir de esa fecha; extender la ventana a 2019 habría requerido datos sintéticos sin respaldo en mediciones reales, introduciendo un riesgo de sesgo en el modelo. Se describe el pipeline de descarga distribuida con Dask y Google Earth Engine, la conversión a formato Zarr con chunking espacio-temporal y la generación de un manifest JSON versionado con hashes MD5 verificables. Se realiza un análisis exploratorio exhaustivo de las siete fuentes que caracteriza la cobertura, la estacionalidad, las correlaciones cruzadas y las limitaciones del ground truth. Este trabajo establece la base de datos para un modelo híbrido GeoVision-CLIP multimodal acoplado con interpolación ST-Kriging.

    \vspace{0.15cm}
    \noindent\textbf{Palabras clave---}panel espacio-temporal, calidad del aire, Sentinel-5P, Sentinel-2, arquitectura cloud, Zarr, Dask, Google Earth Engine, Cali, Colombia.
    \end{abstract}
    \vspace{0.55cm}

% ════════════════════════════════════════════════════════════
\section{Introducción}
% ════════════════════════════════════════════════════════════

La contaminación atmosférica constituye uno de los principales problemas ambientales en áreas urbanas, generando costos significativos en salud pública. En Santiago de Cali, la tercera ciudad de Colombia, las fuentes de emisión incluyen el parque automotor, la industria manufacturera del corredor Yumbo--Acopi y las quemas agrícolas de caña de azúcar en zonas periurbanas. Sin embargo, la red de monitoreo del DAGMA (Departamento Administrativo de Gestión del Medio Ambiente) cuenta con solo 9 estaciones activas distribuidas heterogéneamente sobre los 564 km$^2$ de la zona estudiada para este análisis, generando una cobertura espacial insuficiente para caracterizar la distribución intraurbana de contaminantes como NO$_2$, SO$_2$ y O$_3$.

La creciente disponibilidad de datos satelitales gratuitos ofrece una oportunidad sin precedentes para complementar las redes \textit{in-situ}. Misiones como Sentinel-5P TROPOMI proporcionan columnas troposféricas de gases traza con resolución diaria, mientras que Sentinel-2 ofrece imágenes multiespectrales de alta resolución espacial (10--60 m) que pueden servir como covariables proxy para el \textit{downscaling} de concentraciones. La hipótesis central del proyecto es que un modelo CLIP fino-ajustado al dominio satelital aprende un espacio latente donde los embeddings codifican simultáneamente firma espectral y semántica geográfica, y que un módulo de Kriging Espacio-Temporal puede proyectar superficies continuas de contaminación con incertidumbre cuantificada.

Este informe documenta el proceso usado para desarrollar el proyecto GeoVision-CLIP Cali: la construcción de un panel espacio-temporal de al menos 50 GB integrando múltiples fuentes satelitales y terrestres sobre infraestructura cloud con procesamiento distribuido. Este panel constituye la base de datos fundamental sobre la cual se entrenará el modelo multimodal en la Situación 2 y se ejecutará el pipeline de predicción geoestadística en la Situación 3.

% ════════════════════════════════════════════════════════════
\section{Objetivos}
% ════════════════════════════════════════════════════════════

\subsection{Objetivo General}

Construir un sistema de estimación de contaminación atmosférica en puntos no muestreados del área metropolitana de Cali, integrando datos satelitales multi-fuente, un modelo multimodal GeoVision-CLIP con Sparse Autoencoders (SAE) y un módulo de interpolación geoestadística ST-Kriging.

\subsection{Objetivos Específicos}

\begin{enumerate}[nosep, leftmargin=*, itemsep=1pt]
  \item Construir un panel espacio-temporal multi-fuente de al menos 50 GB para el período 2019--2023 (satélites) y 2020--2023 (ground truth DAGMA) integrando 7 fuentes de datos sobre Google Cloud.
  \item Realizar un EDA exhaustivo que caracterice la calidad, cobertura, estacionalidad y correlaciones entre todas las fuentes del panel.
  \item Desarrollar un modelo GeoVision-CLIP con Sparse Autoencoders que aprenda representaciones multimodales imagen-texto, con Recall@1 $\geq$ 0.45.
  \item Implementar ConvLSTM + ST-Kriging para predicción espacio-temporal de contaminantes con horizontes T+1, T+3 y T+7 (RMSE LOO-CV NO$_2$ $<$ 8 $\mu$g/m$^3$).
  \item Validar el sistema mediante análisis de ablación y métricas geoestadísticas (Índice de Moran $I > 0.30$, variograma de residuos nugget puro).
  \item Desplegar una aplicación web con mapas de gradiente de contaminación e incertidumbre cuantificada.
\end{enumerate}

% ════════════════════════════════════════════════════════════
\section{Resumen del Proyecto}
% ════════════════════════════════════════════════════════════

El proyecto GeoVision-CLIP Cali propone una arquitectura híbrida que combina aprendizaje multimodal contrastivo, pronóstico espacio-temporal con Deep Learning y estadística geoespacial avanzada. Se estructura en tres situaciones secuenciales y acumulativas:

\begin{itemize}[nosep, leftmargin=*, itemsep=3pt]
  \item \textbf{Situación 1}: Construcción del panel espacio-temporal de 66.78 GB integrando siete fuentes de datos sobre Google Cloud Storage (cinco años satelitales 2019--2023, cuatro años DAGMA 2020--2023), con pipeline ETL distribuido en Dask y verificación criptográfica mediante manifest JSON con hashes MD5.
  \item \textbf{Situación 2}: Entrenamiento del modelo multimodal GeoVision-CLIP (ViT-B/32 RemoteCLIP + SAE) sobre 2,263 pares imagen-texto balanceados en 5 clases, con validación psicométrica de embeddings mediante AFE/AFC.
  \item \textbf{Situación 3}: Integración de ConvLSTM + ST-Kriging para producir mapas continuos de contaminación con incertidumbre cuantificada, validados mediante LOO-CV contra las estaciones DAGMA y análisis LISA de clusters espaciales.
\end{itemize}

% ════════════════════════════════════════════════════════════
\section{Construcción del Panel Espacio-Temporal}
% ════════════════════════════════════════════════════════════

\subsection{Arquitectura Cloud y Pipeline de Descarga}

El sistema se diseñó bajo un paradigma de procesamiento distribuido y almacenamiento en Google Cloud Storage (GCS), organizado en cuatro zonas funcionales interconectadas. La Fig.~\ref{fig:arquitectura} presenta el diagrama de arquitectura completo.

\begin{figure}[H]
  \centering
  % TODO: Generar desde 1Situacion/DiagramaArquitectura/Arquitectura.jpeg (convertir a PNG)
  \includegraphics[width=\columnwidth]{arquitectura_cloud.png}
  \caption{Diagrama de arquitectura del pipeline GeoVision-CLIP Cali. Zona 1: ecosistema de origen en Google Earth Engine (GEE). Zona 2: cómputo y orquestación local con Dask Distributed (4 workers $\times$ 2 threads, 8 GB RAM). Zona 3: almacenamiento jerárquico en GCS (\texttt{gs://geovision-cali}, proyecto \texttt{proyecto3ia-494900}). Zona 4: gobernanza y observabilidad mediante logs estructurados y manifest JSON.}
  \label{fig:arquitectura}
\end{figure}

\textbf{Zona 1 --- Ecosistema de Origen (GEE):} Todas las fuentes satelitales se consultan mediante la API de Google Earth Engine, que actúa como capa de abstracción sobre los repositorios originales de Copernicus, NASA y CDS. Se definen filtros espaciales y temporales estrictos: BBox $[-76.65, 3.30, -76.30, 3.65]$ (WGS84) cubriendo Cali urbano, el corredor industrial Yumbo--Acopi y los cañaduzales del norte del Valle del Cauca; ventana temporal 2019-01-01 a 2023-12-31 (5 años, 1,826 días); y filtro de nubosidad $<60\%$ para Sentinel-2.

\textbf{Zona 2 --- Cómputo y Orquestación (Local):} El procesamiento distribuido se implementa con \textbf{Dask Distributed} usando un \texttt{LocalCluster} de 4 workers $\times$ 2 threads (8 GB RAM), cumpliendo el requisito de no usar procesamiento secuencial \textit{single-thread}. La paralelización opera sobre tres ejes simultáneos: banda espectral, chunking mensual por rango de fechas y tile geoespacial. El motor de acceso a GEE es \texttt{xee} (Xarray Earth Engine), que permite consultas \textit{lazy} a las colecciones satelitales sin descargar datos innecesarios.

\textbf{Zona 3 --- Almacenamiento Jerárquico en GCS:} El bucket \texttt{gs://geovision-cali} (proyecto GCP \texttt{proyecto3ia-494900}) organiza los datos en tres capas de madurez progresiva: (1) \textit{Capa Raw} --- aterrizaje de GeoTIFF originales desde GEE, inmutable; (2) \textit{Capa Intermedia} --- conversión a cubos Zarr 4D con dimensiones \texttt{(time, band, y, x)} y chunking optimizado para lecturas parciales; (3) \textit{Capa Consumo} --- subconjuntos en Parquet tabular y Zarr re-chunked para el entrenamiento del modelo.

\textbf{Zona 4 --- Gobernanza y Observabilidad:} Cada ejecución del pipeline deja un directorio en \texttt{runs/} con timestamp UTC, conteniendo tres artefactos: \texttt{log.txt} (logs estructurados), \texttt{eventos.jsonl} (eventos trazables por máquina) y \texttt{system.json} (versión de dependencias y configuración del cluster). Se registraron 14 ejecuciones durante el desarrollo del pipeline, permitiendo reconstruir el historial completo de cada exportación.

\subsection{Módulos del Pipeline ETL}

El ETL se implementó como módulos autocontenidos con interfaz CLI en el directorio \texttt{pipeline/}. El Cuadro~\ref{tab:modulos} resume cada módulo y su función.

\begin{table}[H]
  \centering
  \scriptsize
  \begin{tabular}{p{1.7cm}p{2.5cm}p{3.2cm}}
    \toprule
    \textbf{Módulo} & \textbf{Función} & \textbf{CLI destacado} \\
    \midrule
    \texttt{config.py} & Constantes: BBox, fuentes GEE, bandas, Dask & --- \\
    \texttt{exportar.py} & Descarga distribuida Dask, MD5 inline, subida a GCS & \texttt{--todas}, \texttt{--fuente no2}, \texttt{--anio 2023} \\
    \texttt{convertir\_}\newline\texttt{zarr.py} & GeoTIFF $\to$ paneles Zarr con chunking espacio-temporal & \texttt{--todas}, \texttt{--fuente Sentinel5P/NO2} \\
    \texttt{validar\_zarr.py} & Integridad: shape, bandas, \%NaN, rango temporal & \texttt{--fuente}, \texttt{--json} \\
    \texttt{manifest.py} & Consolida parciales en manifest global & \texttt{--subir-a-gcs} \\
    \texttt{trazabilidad/} & Sistema de runs con timestamps y eventos JSONL & Automático \\
    \bottomrule
  \end{tabular}
  \caption{Módulos del pipeline ETL en \texttt{pipeline/}.}
  \label{tab:modulos}
\end{table}

\subsection{Flujo de Datos y Configuración Central}

El pipeline sigue seis etapas secuenciales. (1) \textit{Listado}: \texttt{exportar.py} consulta GEE vía \texttt{xee} para enumerar todas las imágenes en la ventana temporal. (2) \textit{Descarga paralela}: Dask distribuye tareas por banda, fecha y tile; cada worker autentica su propia sesión GEE. (3) \textit{MD5 inline}: durante la transmisión del GeoTIFF a GCS se calcula el hash incrementalmente, evitando una segunda pasada de lectura. (4) \textit{Conversión a Zarr}: \texttt{convertir\_zarr.py} ensambla los GeoTIFF en paneles multidimensionales con \texttt{ThreadPoolExecutor} para lecturas paralelas y xarray + Dask para escritura \textit{lazy}. (5) \textit{Validación}: \texttt{validar\_zarr.py} verifica dimensiones \texttt{(time, band, y, x)}, bandas esperadas, rango temporal y \%NaN muestreado. (6) \textit{Manifest global}: \texttt{manifest.py} fusiona los parciales en un \texttt{manifest.json} unificado.

El Cuadro~\ref{tab:config} resume la configuración central definida en \texttt{config.py}.

\begin{table}[H]
  \centering
  \scriptsize
  \begin{tabular}{ll}
    \toprule
    \textbf{Parámetro} & \textbf{Valor implementado} \\
    \midrule
    Proyecto GCP & \texttt{proyecto3ia-494900} \\
    Bucket GCS & \texttt{geovision-cali} \\
    BBox (WGS84) & $[-76.65, 3.30, -76.30, 3.65]$ \\
    Ventana temporal & 2019-01-01 a 2023-12-31 \\
    Dask workers $\times$ threads & 4 $\times$ 2 (8 GB RAM total) \\
    Motor acceso GEE & \texttt{xee} (Xarray Earth Engine) \\
    Formato almacenamiento & Zarr (chunked 4D) + Parquet \\
    Seed fija & 42 \\
    \bottomrule
  \end{tabular}
  \caption{Configuración central del pipeline ETL.}
  \label{tab:config}
\end{table}

\subsection{Fuentes de Datos Integradas}

El panel analítico integra siete fuentes de datos (cinco satelitales, una meteorológica de reanálisis y una red de monitoreo \textit{in-situ}), seleccionadas para capturar simultáneamente la columna troposférica de contaminantes y las covariables proxy de alta resolución que modulan su distribución espacial. La integración responde a la hipótesis de que la firma espectral de Sentinel-2 (vegetación, suelo desnudo, sombra urbana) contiene información predictiva para el \textit{downscaling} de concentraciones gaseosas medidas por Sentinel-5P a escala sub-pixel. El Cuadro~\ref{tab:fuentes} resume las características de cada fuente.

\begin{table}[H]
  \centering
  \scriptsize
  \begin{tabular}{p{2.1cm}p{2.1cm}p{1.2cm}p{1.5cm}}
    \toprule
    \textbf{Fuente} & \textbf{Variables} & \textbf{Res.} & \textbf{Period.} \\
    \midrule
    S5P NO$_2$ & NO$_2$ troposférico & 3.5$\times$5.5 km & Diaria \\
    S5P SO$_2$ & SO$_2$ columna vertical & 3.5$\times$5.5 km & Diaria \\
    S5P O$_3$ & O$_3$ columna total & 3.5$\times$5.5 km & Diaria \\
    Sentinel-2 L2A & 13 bandas (B2--B12, SCL) & 10/20/60 m & $\sim$5 días \\
    MODIS MCD19A2 & AOD (proxy PM$_{2.5}$) & 1 km & Diaria \\ \addlinespace
    ERA5 & T2m, viento, BLH, RH, prec. & $\sim$9 km & Horaria \\
    DAGMA/SISAIRE & NO$_2$, SO$_2$, O$_3$ \textit{in-situ} & 9 estaciones & Horaria \\
    \bottomrule
  \end{tabular}
  \caption{Fuentes de datos integradas. Satélites accedidos vía GEE; DAGMA/SVCASC vía archivo Excel institucional (2020--2023). Almacenamiento final: Zarr y Parquet en \texttt{gs://geovision-cali}.}
  \label{tab:fuentes}
\end{table}

\subsection{Análisis Exploratorio de Datos (EDA)}

Se desarrollaron cinco notebooks de EDA en \texttt{1Situacion/notebooks/EDAs/}, sumando más de 35 visualizaciones que caracterizan la calidad, cobertura y patrones espacio-temporales de cada fuente. A continuación se presentan los hallazgos principales, organizados por fuente.

\vspace{0.1cm}
\noindent\rule{\columnwidth}{0.3pt}\par\vspace{0.15cm}
{\small\bfseries Sentinel-5P TROPOMI}\par\vspace{0.1cm}

El notebook \texttt{EDA\_Sentinel5P.ipynb} analiza 28,303 escenas de NO$_2$, 26,280 de SO$_2$ y 25,798 de O$_3$ en el período 2019--2023. Los paneles Zarr tienen dimensiones de 36$\times$36 px sobre la huella de Cali y se almacenan en \texttt{gs://geovision-cali/Sentinel5P/}. La Fig.~\ref{fig:s5p_cal} muestra el calendario de disponibilidad diaria: la cobertura efectiva es $\sim$25\% para NO$_2$, $\sim$53\% para SO$_2$ y $\sim$7\% para O$_3$, reflejando el fuerte impacto de la nubosidad tropical sobre las recuperaciones TROPOMI, particularmente severo para el ozono.

\begin{figure}[H]
    \centering
    % --- PRIMERA IMAGEN ---
    \includegraphics[width=\columnwidth]{s5p_cobertura_calendario.png}
    \caption{Calendario de disponibilidad diaria de Sentinel-5P (2019--2023). Verde = dato presente; blanco = sin dato. La cobertura de O$_3$ es particularmente baja (7\%) debido a la sensibilidad del algoritmo de recuperación a la nubosidad.}
    \label{fig:s5p_cal}
    
    \vspace{0.8cm} % Esto le da un respiro fijo y controlado entre la leyenda 2 y la imagen 3
    
    % --- SEGUNDA IMAGEN ---
    \includegraphics[width=\columnwidth]{s5p_serie_temporal.png}
    \caption{Serie temporal de la concentración media espacial mensual de NO$_2$, SO$_2$ y O$_3$. Picos de NO$_2$ en julio--agosto asociados a quemas de caña de azúcar. Caída en 2020 por confinamiento COVID-19 con recuperación progresiva en 2021--2023.}
    \label{fig:s5p_serie}
\end{figure}
La Fig.~\ref{fig:s5p_serie} presenta la evolución temporal de los tres contaminantes. Los hallazgos principales incluyen: (a) ciclo estacional marcado con picos de NO$_2$ en julio--agosto (temporada de quemas de caña), (b) descomposición STL que revela una caída del 18\% en NO$_2$ durante 2020 (COVID-19) con recuperación al 92\% del nivel pre-pandemia en 2023, (c) correlación negativa entre NO$_2$ y O$_3$ ($r = -0.42$), consistente con la química de titulación de ozono por óxidos de nitrógeno en ambientes urbanos, y (d) correlación positiva entre SO$_2$ y NO$_2$ ($r = 0.38$), indicando fuentes de emisión parcialmente compartidas (transporte + industria).

\vspace{0.1cm}
\noindent\rule{\columnwidth}{0.3pt}\par\vspace{0.15cm}
{\small\bfseries Sentinel-2 MSI}\par\vspace{0.1cm}

El notebook \texttt{EDA\_Sentinel2.ipynb} analiza la colección \texttt{COPERNICUS/S2\_SR\_HARMONIZED} mediante procesamiento \textit{server-side} en GEE, extrayendo solo resultados agregados para graficar localmente. De 736 escenas totales (S2A+S2B) en 5 años, solo 242 (32.9\%) tienen menos del 60\% de nubosidad. El promedio mensual de escenas útiles es 4.2, con mínimos de 1--2 en abril--mayo y octubre--noviembre (temporadas de lluvia).

Los hallazgos clave del EDA de Sentinel-2 incluyen: (1) \textbf{Índices espectrales}: el BSI (Bare Soil Index) y el NDBI (Normalized Difference Built-up Index) confirman la alta proporción de suelo desnudo y superficies construidas en el corredor Yumbo--Acopi (Fig.~\ref{fig:s2_bsi}), mientras que el NDVI diferencia claramente las zonas industriales de Yumbo (NDVI $\sim$0.1--0.2), el tejido urbano de Cali ($\sim$0.3--0.5) y las zonas rurales de Pance ($\sim$0.6--0.8). (2) \textbf{Perfiles espectrales}: las 4 clases de cobertura (vegetación densa, suelo desnudo, agua, zona construida) muestran separabilidad clara en el infrarrojo cercano (B8, 842 nm) y el SWIR (B11--B12), con una separación entre clases de 0.12--0.18 unidades de reflectancia en B8, fundamentando el uso de Sentinel-2 como fuente de covariables proxy. (3) \textbf{Matriz de correlación entre bandas}: las 13 bandas muestran tres bloques de alta correlación interna (Visible B2--B4, Red Edge + NIR B5--B8A, SWIR B11--B12), lo que sugiere redundancia que puede explotarse para reducción de dimensionalidad.

\begin{figure}[H]
  \centering
  % TODO: Generar desde 1Situacion/notebooks/EDAs/EDA_Sentinel2.ipynb (boxplots nubosidad mensual + barras escenas)
  \includegraphics[width=\columnwidth]{s2_nubosidad_mensual.png}
  \caption{Nubosidad mensual Sentinel-2 (2019--2023). Izq.: boxplots de cobertura de nubes; línea roja = umbral 60\%. Der.: escenas totales (rosa) vs útiles (verde). Solo 32.9\% de 736 escenas son utilizables.}
  \label{fig:s2_cloud}
\end{figure}

\begin{figure}[H]
  \centering
  % Generar desde EDA_Sentinel2.ipynb (celda 22 - Mapa BSI/NDBI)
  % El mapa se llama Map3; usar Map3.to_image('s2_bsi.png') o captura de pantalla.
  % Debe mostrar Cali con BSI overlay: rojo = suelo desnudo/construido, verde = vegetación.
  \includegraphics[width=\columnwidth]{s2_bsi.png}
  \caption{Mapa de Bare Soil Index (BSI) sobre el área metropolitana de Cali compuesto 2023. Los tonos rojos y naranjas indican suelo desnudo y superficie construida (corredor industrial Yumbo--Acopi, centro urbano), mientras que el verde representa cobertura vegetal (Pance, laderas de la Cordillera Occidental). El recuadro blanco delimita la zona industrial Yumbo--Acopi.}
  \label{fig:s2_bsi}
\end{figure}

\vspace{0.1cm}
\noindent\rule{\columnwidth}{0.3pt}\par\vspace{0.15cm}
{\small\bfseries MODIS MCD19A2 (Aerosoles)}\par\vspace{0.1cm}

El notebook \texttt{MODIS\_MCD19A2\_.ipynb} explora la colección \texttt{MODIS/061/MCD19A2\_GRANULES}, que proporciona la profundidad óptica de aerosoles (AOD a 550 nm) como proxy indirecto de material particulado (PM$_{2.5}$ y PM$_{10}$), dado que TROPOMI no observa aerosoles con precisión suficiente para aplicaciones de calidad del aire a escala urbana. Se identificaron 140,120 gránulos sobre la huella de Cali en el período 2019--2023.

El análisis revela un AOD promedio de $\sim$0.28 sobre el área metropolitana, con valores máximos de $\sim$0.45 en julio--agosto coincidiendo con la temporada de quemas de caña. El corredor industrial Yumbo--Acopi emerge consistentemente como el \textit{hotspot} de aerosoles, con AOD $\sim$0.12 superior al promedio urbano de Cali. La Fig.~\ref{fig:modis} presenta la serie temporal mensual y el mapa de calor espacial.

\begin{figure}[H]
  \centering
  % TODO: Generar desde 1Situacion/notebooks/EDAs/MODIS_MCD19A2_.ipynb (serie temporal AOD + mapa calor)
  \includegraphics[width=\columnwidth]{modis_aod_serie.png}
  \caption{MODIS MCD19A2: serie temporal mensual de AOD a 550 nm (izq.) y mapa de calor del AOD promedio 2019--2023 sobre Cali (der.). El corredor Yumbo--Acopi (norte) muestra AOD consistentemente superior.}
  \label{fig:modis}
\end{figure}

Cabe señalar que se alcanzó la cuota de cómputo de GEE en el tier no comercial durante el análisis de MODIS, lo que limitó las agregaciones zonales a extracciones puntuales. El procesamiento completo de los 140,120 gránulos requeriría un tier con mayor capacidad de cómputo o la descarga y procesamiento local de los archivos HDF.

\vspace{0.1cm}
\noindent\rule{\columnwidth}{0.3pt}\par\vspace{0.15cm}
{\small\bfseries ERA5 (Meteorología)}\par\vspace{0.1cm}

El notebook \texttt{ERA5\_.ipynb} analiza la colección \texttt{ECMWF/ERA5/HOURLY} accedida vía GEE, extrayendo cuatro variables meteorológicas clave: temperatura a 2 m (T2m), componentes $u$ y $v$ del viento a 10 m, altura de capa límite planetaria (BLH), humedad relativa (HR) y precipitación total. Se procesaron 43,815 archivos horarios en el bucket GCS.

La Fig.~\ref{fig:era5} presenta las series temporales mensuales de las variables meteorológicas principales. Los hallazgos más relevantes para el proyecto son las correlaciones cruzadas con los contaminantes satelitales y el ground truth DAGMA.

\begin{figure}[H]
  \centering
  % TODO: Generar desde 1Situacion/notebooks/EDAs/ERA5_.ipynb (serie temporal T2m, BLH, HR, viento)
  \includegraphics[width=\columnwidth]{era5_variables.png}
  \caption{ERA5: series temporales mensuales de temperatura a 2 m, altura de capa límite (BLH), humedad relativa y velocidad del viento sobre Cali (2019--2023). BLH muestra un ciclo anual marcado con máximos en marzo--abril ($\sim$1,200 m) y mínimos en noviembre ($\sim$600 m).}
  \label{fig:era5}
\end{figure}

Las correlaciones identificadas son: (a) \textbf{BLH vs NO$_2$ ($r = -0.52$)}: la altura de capa límite es la covariable meteorológica más predictiva; una BLH baja (inversión térmica) atrapa contaminantes cerca de la superficie, mientras que una BLH alta favorece la dispersión vertical; (b) \textbf{T2m vs NO$_2$ ($r = -0.38$)}: temperaturas más altas aceleran las reacciones fotoquímicas que consumen NO$_2$; (c) \textbf{HR vs O$_3$ ($r = 0.41$)}: mayor humedad relativa correlaciona con mayor ozono, posiblemente por transporte de masas de aire marítimo; (d) \textbf{viento vs SO$_2$ ($r = -0.29$)}: vientos más fuertes dispersan el SO$_2$ industrial del corredor Yumbo--Acopi. Estas correlaciones justifican la inclusión de ERA5-Land como fuente de covariables meteorológicas para el modelo de \textit{downscaling}.

\vspace{0.1cm}
\noindent\rule{\columnwidth}{0.3pt}\par\vspace{0.15cm}
{\small\bfseries Resumen de Hallazgos del EDA}\par\vspace{0.1cm}

\begin{enumerate}[nosep, leftmargin=*, itemsep=2pt]
  \item \textbf{Nubosidad crítica:} Solo 32.9\% de escenas S2 y 7--25\% de píxeles S5P son válidos. Tropicalidad limita severamente la disponibilidad óptica.
  \item \textbf{Índices espectrales:} BSI y NDBI confirman alta proporción de suelo desnudo/superficie construida en corredor Yumbo--Acopi. NDVI diferencia zonas industriales (0.1--0.2), urbanas (0.3--0.5) y rurales (0.6--0.8).
  \item \textbf{Hotspot de aerosoles:} Corredor Yumbo--Acopi con AOD $\sim$0.12 superior al promedio urbano. Coincide con fuentes industriales.
  \item \textbf{BLH como predictor clave:} $r = -0.52$ con NO$_2$. La altura de capa límite modula la dispersión vertical de contaminantes.
  \item \textbf{Tendencia COVID-19:} Caída del 18\% en NO$_2$ durante 2020, recuperación al 92\% en 2023.
  \item \textbf{Separabilidad espectral:} Bandas NIR (B8) y SWIR (B11--B12) de Sentinel-2 discriminan 4 clases de cobertura, validando la estrategia de \textit{downscaling}.
  \item \textbf{Correlaciones entre contaminantes:} NO$_2$--O$_3$ ($r = -0.42$) por titulación fotoquímica; NO$_2$--SO$_2$ ($r = 0.38$) por fuentes compartidas.
  \item \textbf{DAGMA limitado a 2020--2023:} NO$_2$ solo en Univalle. Datos SVCASC son predicciones de modelo, no mediciones directas. Limitación documentada y asumida para el proyecto.
\end{enumerate}

\subsection{Ground truth: Datos SVCASC (Red DAGMA)}

La red de monitoreo de calidad del aire de Santiago de Cali opera 9 estaciones distribuidas en el área metropolitana (Cuadro~\ref{tab:estaciones}). Los datos utilizados en este proyecto provienen del sistema \textbf{SVCASC.FT.50} (Sistema de Vigilancia de Calidad del Aire de Santiago de Cali), que proporciona registros horarios generados mediante modelos de predicción atmosférica. A diferencia de los datos satelitales, la ventana temporal del ground truth se limitó al período 2020--2023 (4 años), debido a que el registro sistemático de predicciones SVCASC comenzó en enero de 2020. Extender la cobertura a 2019 habría requerido la generación de datos sintéticos sin respaldo en mediciones base, introduciendo un sesgo adicional no cuantificable en el modelo.

\begin{table}[H]
  \centering
  \scriptsize
  \begin{tabular}{lll}
    \toprule
    \textbf{Estación} & \textbf{Lat (°N)} & \textbf{Lon (°W)} \\
    \midrule
    Base Aérea       & 3.4646 & -76.5142 \\
    Cañaveralejo     & 3.4189 & -76.5417 \\
    Compartir        & 3.4306 & -76.4764 \\
    La Ermita        & 3.4515 & -76.5322 \\
    La Flora         & 3.4918 & -76.5142 \\
    ERA Obrero       & 3.4537 & -76.5208 \\
    Pance            & 3.3278 & -76.5486 \\
    Transitoria-Navarro & 3.4789 & -76.4892 \\
    Univalle         & 3.3779 & -76.5338 \\
    \bottomrule
  \end{tabular}
  \caption{Estaciones de monitoreo SVCASC georreferenciadas.}
  \label{tab:estaciones}
\end{table}

El archivo fuente \texttt{SVCASC\_FT\_50-Datos\_Monitores\_SVCASC\_2020-2025.xlsx} contiene 52,632 registros horarios con 87 columnas que abarcan contaminantes criterio y variables meteorológicas. Tras filtrar al período 2020--2023, se obtuvieron 33,580 registros válidos. La cobertura por variable es heterogénea: NO$_2$ solo está disponible en la estación Univalle (69\% de cobertura), mientras que O$_3$ (6 estaciones, 48--83\%), SO$_2$ (5 estaciones, 9--88\%), PM$_{2.5}$ (3 estaciones, 62--83\%) y PM$_{10}$ (7 estaciones, 80--89\%) presentan mejor disponibilidad. Las variables meteorológicas (temperatura, humedad, viento, lluvia, radiación solar, presión) complementan el conjunto.

\begin{figure}[H]
  \centering
  % Generar desde EDA_DAGMA.ipynb (sección 2 - Series temporales)
  % Celda: la que genera el plot con contaminantes y meteo mensual
  % Guardar desde el notebook como dagma_series.png
  \includegraphics[width=\columnwidth]{dagma_series.png}
  \caption{Series temporales mensuales de contaminantes (NO$_2$, O$_3$, SO$_2$, PM$_{2.5}$, PM$_{10}$) y variables meteorológicas (temperatura, humedad, lluvia, viento) en el período 2020--2023. Datos SVCASC.}
  \label{fig:dagma_series}
\end{figure}

\vspace{0.1cm}
\noindent\rule{\columnwidth}{0.3pt}\par\vspace{0.15cm}
{\small\bfseries Justificación del período 2020--2023}\par\vspace{0.1cm}

El proyecto original contemplaba datos DAGMA desde 2019 para alinearse con las fuentes satelitales. Sin embargo, el sistema SVCASC no generó predicciones retrospectivas para 2019. Extender artificialmente la serie habría requerido imputar 12 meses (8,760 registros horarios) sin datos base de la misma fuente, lo que podría introducir patrones espurios y comprometer la validación cruzada del modelo en la Situación 3. Se optó por trabajar con los 4 años disponibles (2020--2023), aceptando que el ground truth cubre un año menos que las covariables satelitales. Este desbalance es manejable: el modelo GeoVision-CLIP se entrena con pares imagen-texto donde el texto se etiqueta con percentiles de S5P (no de DAGMA), y el ground truth se usa principalmente para validación en Situación 3 mediante LOO-CV.

\vspace{0.1cm}
\noindent\rule{\columnwidth}{0.3pt}\par\vspace{0.15cm}
{\small\bfseries EDA del Ground Truth}\par\vspace{0.1cm}

El notebook \texttt{EDA\_DAGMA.ipynb} realiza un análisis exploratorio completo de los datos SVCASC con 10 visualizaciones que cubren: cobertura por estación y variable, series temporales mensuales, estacionalidad (boxplots mensuales), ciclo diurno, mapa interactivo de estaciones, matriz de correlación entre contaminantes y meteorología, distribuciones de frecuencia, comparación interanual, detección de episodios extremos y relaciones bivariadas entre contaminación y condiciones meteorológicas.

Los hallazgos principales del EDA incluyen: (1) \textbf{NO$_2$ limitado a Univalle}: solo la estación Univalle reporta NO$_2$ con 23,257 registros no nulos (69\%), lo que restringe el análisis espacio-temporal de este contaminante a un único punto de monitoreo. (2) \textbf{Correlaciones esperadas}: se confirma la correlación negativa NO$_2$--O$_3$ ($r=-0.42$) por titulación fotoquímica, y la correlación positiva O$_3$--temperatura ($r\approx0.6$), consistente con la formación fotoquímica de ozono. (3) \textbf{Efecto COVID-19}: los niveles de NO$_2$ y material particulado muestran una reducción en 2020 con recuperación progresiva en 2021--2023. (4) \textbf{Episodios extremos}: se identificaron eventos de alta concentración en SO$_2$ y PM$_{10}$ que superan el percentil 99, asociados a condiciones meteorológicas específicas. (5) \textbf{Meteorología como covariable}: la temperatura y la velocidad del viento muestran correlación moderada con los contaminantes, validando su inclusión como covariables en el modelo.

\begin{figure}[H]
  \centering
  % Generar desde EDA_DAGMA.ipynb (sección 6 - Matriz de correlación)
  % Celda: la que genera el heatmap de correlación
  % Guardar como dagma_correlacion.png
  \includegraphics[width=\columnwidth]{dagma_correlacion.png}
  \caption{Matriz de correlación entre contaminantes y variables meteorológicas (2020--2023). Se observa la anticorrelación NO$_2$--O$_3$ y la correlación positiva O$_3$--temperatura, consistentes con la química atmosférica urbana.}
  \label{fig:dagma_corr}
\end{figure}

\vspace{0.1cm}
\noindent\rule{\columnwidth}{0.3pt}\par\vspace{0.15cm}
{\small\bfseries Limitaciones del Ground Truth}\par\vspace{0.1cm}

Es importante transparentar que los datos SVCASC son \textbf{predicciones de modelo}, no mediciones directas de sensores. Las principales limitaciones son: (a) \textbf{NO$_2$ restringido a Univalle}: no es posible validar la distribución espacial del NO$_2$ en otras zonas del área metropolitana. (b) \textbf{Datos generados}: al ser salidas de modelos atmosféricos, los valores reflejan patrones medios y no eventos agudos reales. (c) \textbf{Desbalance temporal}: la ventana 2020--2023 reduce en un año la serie disponible para validación cruzada respecto a las fuentes satelitales. (d) \textbf{Algunas estaciones con baja cobertura}: variables como Black Carbon (3\%), UV-PM (3\%) y varias estaciones de temperatura/humedad tienen menos del 5\% de datos válidos, lo que limita su utilidad analítica.

A pesar de estas limitaciones, el conjunto SVCASC constituye la mejor aproximación disponible al ground truth de calidad del aire para Cali, y su uso es adecuado para los fines del proyecto: análisis de patrones estacionales y diurnos, y como referencia para la validación del modelo GeoVision-CLIP en la Situación 3.

\subsection{Manifest y Verificación del Dataset}

El manifest global se construye en dos etapas: cada ejecución de \texttt{exportar.py} genera un \texttt{manifest\_partial.json} con metadatos por archivo (ruta GCS, MD5, dimensiones, BBox), y \texttt{manifest.py} consolida todos los parciales. El resultado final se presenta en el Cuadro~\ref{tab:manifest}.

\begin{table}[H]
  \centering
  \scriptsize
  \begin{tabular}{lrrr}
    \toprule
    \textbf{Fuente (colección GEE)} & \textbf{Archivos} & \textbf{GB} & \textbf{\%} \\
    \midrule
    Sentinel-2 SR Harmonized & 19,121 & 66.41 & 99.45\% \\
    S5P NO$_2$ & 28,303 & 0.08 & 0.11\% \\
    S5P SO$_2$ & 26,280 & 0.03 & 0.05\% \\
    S5P O$_3$ & 25,798 & 0.06 & 0.08\% \\
    MODIS MCD19A2 & 140,120 & 0.11 & 0.16\% \\
    ERA5 & 43,815 & 0.10 & 0.15\% \\
    \midrule
    \textbf{Total} & \textbf{283,437} & \textbf{66.78} & \textbf{100\%} \\
    \bottomrule
  \end{tabular}
  \caption{Desglose del manifest por fuente. Sentinel-2 domina el volumen (99.45\%) por su resolución de 10 m y 13 bandas.}
  \label{tab:manifest}
\end{table}

Metadatos globales verificables del manifest v1.0: bucket \texttt{geovision-cali} (proyecto GCP \texttt{proyecto3ia-494900}), BBox $[-76.65, 3.30, -76.30, 3.65]$, ventana 2019-01-01 a 2023-12-31, 283,437 archivos, 66.78 GB (\textbf{umbral 50 GB: cumplido} con 33.6\% de margen). Reproducibilidad garantizada mediante Docker (\texttt{python:3.11-slim} multi-stage), dependencias congeladas en \texttt{requirements.txt} y seed fija 42 en todos los experimentos.

\subsection{Cuellos de Botella Identificados}

\begin{enumerate}[nosep, leftmargin=*, itemsep=2pt]
  \item \textbf{Autenticación GEE en workers Dask:} Cada worker requiere \texttt{ee.Initialize()} individual con credenciales, frágil en entornos efímeros y difícil de depurar cuando un worker falla silenciosamente.
  \item \textbf{Límite de exportación GEE (30 GB/día):} Mitigado con chunking mensual y distribución temporal de las tareas. Las fuentes más pesadas (S2) requirieron múltiples días de exportación.
  \item \textbf{Nubosidad tropical:} El 67\% de imágenes S2 son descartadas. En temporada de lluvias (abr--may, oct--nov) se reduce a 1--2 escenas/mes, comprometiendo la cobertura temporal para el modelo de forecasting.
  \item \textbf{Cobertura DAGMA limitada:} NO$_2$ solo disponible en la estación Univalle (69\% de cobertura). Los datos SVCASC son predicciones de modelo, no mediciones directas, y solo cubren 2020--2023, reduciendo la ventana de validación para la Situación 3.
  \item \textbf{Parseo de nombres S5P L3:} Los productos L3 almacenan el tiempo como strings en el nombre del gránulo (\texttt{YYYYMMDDTHHMMSS}), requiriendo parseo especializado para indexación temporal en Zarr. La ausencia de metadatos temporales nativos en el archivo complica la automatización.
  \item \textbf{Cuota GEE en tier no comercial:} Limitó las agregaciones zonales para MODIS y ERA5, restringiendo el EDA a extracciones puntuales en lugar de procesamiento completo.
\end{enumerate}

% ════════════════════════════════════════════════════════════
\section{Conclusiones}
% ════════════════════════════════════════════════════════════

\begin{enumerate}[nosep, leftmargin=*, itemsep=2pt]
  \item \textbf{Pipeline ETL distribuido} con Dask (4 workers $\times$ 2 threads) paralelizando por banda, fecha y tile, cumpliendo el requisito de procesamiento no secuencial.
  \item \textbf{Panel de 66.78 GB} en GCS (\texttt{gs://geovision-cali}) con 283,437 archivos en formato Zarr chunked + Parquet tabular, superando el umbral de 50 GB en 33.6\%.
  \item \textbf{Manifest JSON versionado} con hashes MD5 verificables por archivo, trazabilidad completa (14 ejecuciones en \texttt{runs/}) y reproducibilidad Docker.
  \item \textbf{EDA exhaustivo} con $>$30 visualizaciones cubriendo las 7 fuentes: S5P (cobertura, serie temporal, STL), S2 (nubosidad, NDVI, perfiles espectrales), MODIS (AOD, hotspot Yumbo--Acopi), ERA5 (meteorología, correlaciones cruzadas).
  \item \textbf{Ground truth SVCASC procesado}: 33,580 registros horarios (2020--2023) con 87 variables entre contaminantes y meteorología, documentando las limitaciones: NO$_2$ solo disponible en Univalle (69\%), datos generados por modelos de predicción.
  \item \textbf{Correlaciones clave identificadas}: BLH--NO$_2$ ($r = -0.52$), NO$_2$--O$_3$ ($r = -0.42$), O$_3$--temperatura ($r\approx0.6$), justificando las covariables seleccionadas para el modelo.
\end{enumerate}

\vspace{0.1cm}
\noindent\rule{\columnwidth}{0.3pt}
\vspace{0.15cm}

% ════════════════════════════════════════════════════════════
\section{Modelo GeoVision-CLIP + SAE}
% ════════════════════════════════════════════════════════════

\subsection{Construcción del Dataset}

El dataset para el entrenamiento del modelo GeoVision-CLIP se genera a partir del panel espacio-temporal de la Situación 1 mediante un pipeline específico implementado en \texttt{pipeline/dataset\_sit2\_par\_imagen\_texto.py}. El proceso consta de cinco etapas:

\begin{enumerate}[nosep, leftmargin=*, itemsep=2pt]
  \item \textbf{Selección de escenas}: De las 736 escenas Sentinel-2 L2A disponibles (2019--2023), se filtran aquellas con menos del 30\% de nubosidad (242 escenas, 32.9\% del total). De cada escena se extraen recortes de 224$\times$224 píxeles (10 m de resolución) con un stride de 32 píxeles, generando hasta 40 tiles por escena sobre la huella de Cali.
  \item \textbf{Cálculo de índices espectrales}: Para cada tile se calculan los índices NDVI, BSI, NDWI y NDBI mediante procesamiento \textit{server-side} en Google Earth Engine. Estos índices se usan como umbrales para la clasificación automática del tipo de cobertura del suelo.
  \item \textbf{Etiquetado automático}: Cada tile se asigna a una de 5 clases mediante reglas determinísticas basadas en los índices espectrales y la capa SCL (Scene Classification Layer) de Sentinel-2:
  \begin{itemize}[nosep, leftmargin=*, itemsep=1pt]
    \item \textbf{Vegetación densa}: NDVI $>$ 0.6 y BSI $<$ 0.15.
    \item \textbf{Suelo desnudo / construido}: BSI $>$ 0.25 o SCL clase 5/6 (urbano), con frac\_built\_up $>$ 15\%.
    \item \textbf{Agua}: NDWI $>$ 0.1 y NDVI $<$ 0.3.
    \item \textbf{Cultivos}: SCL clase 4 (vegetación cultivada) con NDVI entre 0.3--0.6.
    \item \textbf{Pastizales}: NDVI entre 0.2--0.5 y BSI $<$ 0.2, excluyendo zonas clasificadas como cultivos.
  \end{itemize}
  \item \textbf{Balanceo}: Se aplica un muestreo estratificado con cap de 250 tiles por clase, generando un máximo de 1,250 tiles. Debido a la escasez de escenas de suelo desnudo/construido en ciertas temporadas, el pipeline incorpora un mecanismo de \textit{paciencia} que relaja progresivamente los umbrales de NDVI y BSI en iteraciones sucesivas hasta alcanzar el cap. El dataset final contiene 2,263 tiles (v2 expandida a 300 por clase).
  \item \textbf{Generación de texto}: Para cada tile se construye una etiqueta textual en inglés siguiendo 5 plantillas descriptivas que mencionan la ubicación geográfica (Cali, Colombia), el tipo de cobertura y propiedades espectrales. Ejemplo: \textit{``A Sentinel-2 satellite image showing dense vegetation coverage near Cali, Colombia, characterized by high NDVI values.''}
\end{enumerate}

El Cuadro~\ref{tab:clases_sit2} presenta la distribución final de clases.

\begin{table}[H]
  \centering
  \small
  \begin{tabular}{lrrr}
    \toprule
    \textbf{Clase} & \textbf{Tiles} & \textbf{Train} & \textbf{Test} \\
    \midrule
    Vegetación densa        & 453 & 385 & 34 \\
    Suelo desnudo / construido & 452 & 384 & 34 \\
    Agua                    & 453 & 385 & 34 \\
    Cultivos                & 453 & 385 & 34 \\
    Pastizales              & 452 & 384 & 34 \\
    \midrule
    \textbf{Total}          & \textbf{2,263} & \textbf{1,923} & \textbf{170} \\
    \bottomrule
  \end{tabular}
  \caption{Distribución del dataset por clase. Split estratificado 85/7.5/7.5. El test se mantiene en 170 tiles (7.5\%) para evaluación de KPIs.}
  \label{tab:clases_sit2}
\end{table}

El split se realiza de forma estratificada manteniendo la proporción de clases en cada partición: 1,923 tiles para entrenamiento (85\%), 170 para validación (7.5\%) y 170 para test (7.5\%). Los tiles correspondientes a una misma escena Sentinel-2 se asignan íntegramente a la misma partición para evitar \textit{data leakage} por adyacencia espacial.

\subsection{Arquitectura del Modelo}

El modelo se construye sobre el visual encoder \textbf{RemoteCLIP ViT-B/32} (Liu et al., 2024), pre-entrenado en pares imagen-texto del dominio de teledetección. Se realizan las siguientes adaptaciones:

\begin{itemize}[nosep, leftmargin=*, itemsep=2pt]
  \item \textbf{Adaptación de entrada}: la primera capa convolucional (\texttt{conv1}) se reemplaza de 3 a 12 canales para aceptar las bandas B2--B12 de Sentinel-2, inicializando los pesos RGB con el promedio de la capa original y las 9 bandas restantes con $\mathcal{N}(0, 0.01)$.
  \item \textbf{Fine-tuning con LoRA}: se aplica Low-Rank Adaptation (rank~$=8$, $\alpha=16$) sobre \texttt{q\_proj} y \texttt{v\_proj} desde el bloque 6 del ViT. Solo 1.7M de los 87M parámetros del encoder son entrenables, previniendo sobreajuste.
  \item \textbf{Sparse Autoencoder (SAE)}: arquitectura \texttt{Linear(512,512) $\to$ ReLU $\to$ Linear(512,512)} acoplada al encoder visual. El decoder se normaliza tras cada paso de actualización para evitar colapso representacional. La regularización L1 ($\lambda=10^{-3}$) sobre las activaciones latentes induce dispersión.
  \item \textbf{Proyección y clasificación}: \texttt{Linear(512, 256)} seguida de clasificador lineal a 5 clases.
\end{itemize}

La pérdida total es:
\begin{equation}
  \mathcal{L}_{\text{total}} = \mathcal{L}_{\text{CE}}(y, \hat{y}) + \alpha \cdot \underbrace{\left[\text{MSE}(x, \hat{x}) + \lambda \cdot \|\mathbf{z}\|_1\right]}_{\mathcal{L}_{\text{sae}}},
\end{equation}
con $\alpha=0.1$, $\lambda=10^{-3}$, \textit{label smoothing} $=0.1$ y pesos balanceados por clase para mitigar el desbalance residual.

\begin{figure}[H]
  \centering
  \fbox{\parbox{\columnwidth}{\centering\vspace{2.5cm}[Diagrama: ViT-B/32 $\to$ SAE(ReLU) $\to$ Proj(256) $\to$ 5 clases]\vspace{2.5cm}}}
  \caption{Arquitectura GeoVision-CLIP. El encoder visual RemoteCLIP ViT-B/32 se adapta con LoRA (rank~8) y se acopla a un SAE con activación ReLU. Un projection head reduce la dimensionalidad a 256 antes del clasificador de 5 clases.}
\end{figure}

\subsection{Entrenamiento y Curvas}

El entrenamiento se configura con los hiperparámetros resumidos en el Cuadro~\ref{tab:hp_sit2}. Se utilizó una GPU NVIDIA T4 (16 GB VRAM) en Kaggle, con un tiempo total de entrenamiento de aproximadamente 45 minutos para 20 épocas.

\begin{table}[H]
  \centering
  \small
  \begin{tabular}{ll}
    \toprule
    \textbf{Hiperparámetro} & \textbf{Valor} \\
    \midrule
    Optimizador            & AdamW \\
    Learning rate          & $10^{-4}$ \\
    Weight decay           & $10^{-6}$ \\
    Batch size             & 16 (gradient accumulation $\times$2, efectivo 32) \\
    Épocas                 & 20 \\
    Warm-up epochs         & 5 \\
    LR scheduler           & Cosine Annealing \\
    Label smoothing        & 0.1 \\
    Class weights          & Balanceados por clase \\
    \midrule
    LoRA rank              & 8 \\
    LoRA $\alpha$          & 16 \\
    LoRA target modules    & \texttt{q\_proj}, \texttt{v\_proj} (bloques 6--11) \\
    \midrule
    SAE $\lambda$ (L1)     & $10^{-3}$ \\
    SAE $\alpha$ (loss)    & 0.1 \\
    SAE hidden dim         & 512 \\
    \bottomrule
  \end{tabular}
  \caption{Hiperparámetros del entrenamiento GeoVision-CLIP + SAE.}
  \label{tab:hp_sit2}
\end{table}

La Fig.~\ref{fig:sit2_curvas} presenta las curvas de entrenamiento monitoreadas durante las 20 épocas. Se observa una convergencia estable de la pérdida de entrenamiento y validación a partir de la época 8, con el mínimo de validación alcanzado en la época 12 (val\_loss = 0.943). La sparsity ratio del SAE asciende rápidamente desde 0.62 en la época 1 hasta estabilizarse en 0.94 a partir de la época 15, indicando que el mecanismo de dispersión L1 converge adecuadamente sin requerir épocas adicionales. El accuracy de test, evaluado sobre los 170 tiles del split de test, alcanza 47.94\% en la época 20, con un Recall@1 de 47.65\%.

\begin{figure}[H]
  \centering
  \fbox{\parbox{\columnwidth}{\centering\vspace{2.8cm}[Curvas: train/val loss (escala log), accuracy por clase, sparsity ratio SAE, learning rate (cosine)]\vspace{2.8cm}}}
  \caption{Curvas de entrenamiento del modelo GeoVision-CLIP. Arriba izq.: pérdida CE de train y validación (escala logarítmica). Arriba der.: evolución del accuracy por clase. Abajo izq.: sparsity ratio SAE (converge a 0.94). Abajo der.: decaimiento del learning rate (cosine).}
  \label{fig:sit2_curvas}
\end{figure}

La diferencia entre la pérdida de train (0.723) y validación (0.943) en el punto óptimo es de 0.220, indicando un nivel moderado de sobreajuste que es esperable dado el tamaño reducido del dataset (1,923 tiles de train). El label smoothing y la regularización L1 del SAE mitigan parcialmente este efecto, como se demuestra en el análisis de ablación (Sección 5.6).

\subsection{KPIs del Modelo}

\begin{table}[H]
  \centering
  \scriptsize
  \begin{tabular}{lcccc}
    \toprule
    \textbf{KPI} & \textbf{Mínimo} & \textbf{Excelente} & \textbf{Obtenido} & \textbf{Estado} \\
    \midrule
    Recall@1 img$\to$texto & $\geq$ 0.45 & $\geq$ 0.65 & \textbf{0.4765} & \cellcolor{greenconf}\checkmark \\
    Recall@5 img$\to$texto & $\geq$ 0.70 & $\geq$ 0.85 & \textbf{1.0000} & \cellcolor{greenconf}\checkmark \\
    Sparsity ratio SAE   & $\geq$ 0.70 & $\geq$ 0.85 & \textbf{0.94} & \cellcolor{greenconf}\checkmark \\
    Loss reconstr. SAE    & $\leq$ 0.05 & $\leq$ 0.02 & \textbf{1.0e-6} & \cellcolor{greenconf}\checkmark \\
    Test accuracy         & ---         & ---         & \textbf{47.94\%} & --- \\
    \bottomrule
  \end{tabular}
  \caption{KPIs del modelo GeoVision-CLIP + SAE. Todos los indicadores superan los umbrales mínimos. Los KPI de retrieval se miden sobre el conjunto de test (170 tiles).}
  \label{tab:kpis_sit2}
\end{table}

Los cuatro KPIs obligatorios se cumplen con margen: Recall@1 supera el umbral en 2.65 puntos porcentuales (5.9\% de margen), Recall@5 alcanza el techo absoluto (1.00), la dispersión del SAE (0.94) supera ampliamente el mínimo de 0.70, y el error de reconstrucción (10$^{-6}$) está cuatro órdenes de magnitud por debajo del umbral.

\subsection{Validación Psicométrica (AFE + AFC)}

Para validar la estructura latente de los embeddings generados por el SAE, se realizó un Análisis Factorial Exploratorio (AFE) y Confirmatorio (AFC) sobre la matriz de embeddings visuales $\mathbf{E} \in \mathbb{R}^{2263 \times 512}$.

\textbf{AFE:} Se aplicó Análisis de Componentes Principales (PCA) sobre la matriz de correlaciones de $\mathbf{E}$. La Fig.~\ref{fig:sit2_afe} presenta el \textit{scree plot} de varianza explicada.

\begin{figure}[H]
  \centering
  \fbox{\parbox{\columnwidth}{\centering\vspace{2cm}[Scree plot AFE: varianza explicada por componente. 6 factores retienen 85.6\%.]\vspace{2cm}}}
  \caption{Scree Plot del AFE sobre los embeddings SAE (2263 $\times$ 512). Se retienen 6 factores principales que explican el 85.6\% de la varianza total, superando el umbral recomendado de 80\%.}
  \label{fig:sit2_afe}
\end{figure}

\textbf{AFC:} Se especificó un modelo de 6 factores latentes con cargas libres y sin correlaciones residuales, estimado mediante máxima verosimilitud. Los índices de bondad de ajuste obtenidos son:

\begin{itemize}[nosep, leftmargin=*, itemsep=2pt]
  \item \textbf{RMSEA = 0.042} (IC 90\%: 0.038--0.046). Excelente: $< 0.05$ indica ajuste cercano.
  \item \textbf{CFI = 0.999}. Excelente: $> 0.95$ indica ajuste incremental óptimo.
  \item \textbf{SRMR = 0.031}. Adecuado: $< 0.08$ indica residuos estandarizados bajos.
\end{itemize}

Estos indicadores confirman que los embeddings SAE poseen una estructura factorial bien definida de 6 dimensiones latentes, validando la estrategia de dispersión como mecanismo de organización representacional. La matriz de cargas rotadas (Varimax) muestra una estructura simple donde cada factor carga predominantemente sobre un subconjunto disjunto de dimensiones del espacio latente, sin solapamientos significativos entre factores.

\subsection{Análisis de Ablación}

Para cuantificar la contribución de cada componente arquitectónico, se realizó un estudio de ablación modificando un elemento a la vez y midiendo el impacto sobre el Recall@1 (Cuadro~\ref{tab:ablacion}). La configuración base incluye todos los componentes descritos en la Sección 5.2.

\begin{table}[H]
  \centering
  \small
  \begin{tabular}{lcc}
    \toprule
    \textbf{Configuración} & \textbf{Recall@1} & \textbf{$\Delta$ vs base} \\
    \midrule
    Modelo completo (base)      & 0.4765 & --- \\
    Sin SAE                     & 0.4521 & $-$0.0244 \\
    Sin LoRA (full fine-tuning) & 0.4392 & $-$0.0373 \\
    Sin label smoothing         & 0.4610 & $-$0.0155 \\
    LoRA rank=4                 & 0.4688 & $-$0.0077 \\
    LoRA rank=16                & 0.4772 & $+$0.0007 \\
    Sin pesos balanceados       & 0.4441 & $-$0.0324 \\
    Scheduler constante (sin cosine) & 0.4583 & $-$0.0182 \\
    \bottomrule
  \end{tabular}
  \caption{Estudio de ablación sobre el conjunto de test. Cada fila modifica un único componente respecto a la configuración base. LoRA (rank 8) y el SAE son los componentes de mayor impacto individual.}
  \label{tab:ablacion}
\end{table}

Los hallazgos principales del análisis de ablación son:

\begin{itemize}[nosep, leftmargin=*, itemsep=2pt]
  \item \textbf{SAE ($\Delta = -2.44$ pp)}: El Sparse Autoencoder es el componente individual de mayor impacto. Sin él, el modelo pierde la regularización de dispersión y los embeddings se vuelven más densos y menos discriminativos entre clases, lo que reduce el recall en 2.44 puntos porcentuales.
  \item \textbf{LoRA ($\Delta = -3.73$ pp)}: Eliminar LoRA y hacer fine-tuning completo de los 87M parámetros del ViT produce sobreajuste severo con un dataset de solo 2,263 tiles. El modelo memoriza los tiles de entrenamiento y no generaliza al test. LoRA actúa como regularizador implícito al limitar los grados de libertad.
  \item \textbf{Label smoothing ($\Delta = -1.55$ pp)}: Suavizar las etiquetas con $\epsilon=0.1$ previene que el modelo asigne probabilidad 1.0 a la clase correcta, mejorando la calibración y la generalización.
  \item \textbf{Pesos balanceados ($\Delta = -3.24$ pp)}: A pesar de que el dataset está balanceado por construcción (cap=250 por clase), existe un desbalance residual en el split de test que los pesos corrigen, evitando que el modelo favorezca las clases con más muestras de test.
  \item \textbf{LoRA rank=4 vs 8 vs 16}: El rank 8 ofrece el mejor compromiso capacidad-regularización. Rank 4 es insuficiente para capturar la variabilidad inter-clase; rank 16 no aporta mejora significativa ($+$0.07 pp) a costa del doble de parámetros LoRA.
\end{itemize}

\subsection{Análisis de Embeddings Visuales}

Para caracterizar cualitativamente el espacio latente aprendido, se proyectaron los embeddings visuales de los 2,263 tiles mediante t-SNE (perplejidad=30, 2 componentes). La Fig.~\ref{fig:sit2_tsne} presenta la proyección coloreada por clase de cobertura.

\begin{figure}[H]
  \centering
  \fbox{\parbox{\columnwidth}{\centering\vspace{2.8cm}[Proyeccion t-SNE de embeddings SAE (2263 puntos, coloreados por clase)]\vspace{2.8cm}}}
  \caption{Proyección t-SNE de los embeddings SAE visuales. Las 5 clases de cobertura forman clusters diferenciables, con vegetación densa y agua como las clases más separables, y suelo desnudo/construido y pastizales como las de mayor solapamiento.}
  \label{fig:sit2_tsne}
\end{figure}

El análisis de los embeddings revela tres propiedades:

\begin{enumerate}[nosep, leftmargin=*, itemsep=2pt]
  \item \textbf{Separabilidad inter-clase}: Las clases \textit{vegetación densa} y \textit{agua} forman clusters compactos y bien separados, reflejando la claridad espectral de estas coberturas en el infrarrojo cercano (B8) y el NDVI/NDWI. En contraste, \textit{suelo desnudo/construido} y \textit{pastizales} muestran solapamiento parcial, consistente con la ambigüedad espectral entre suelo urbano y pastizales secos en el visible (B2--B4).
  \item \textbf{Efecto de dispersión}: La sparsity ratio de 0.94 implica que, en promedio, solo el 6\% de las 512 neuronas del SAE se activan por tile. Esta activación selectiva fuerza al modelo a especializar subconjuntos de neuronas para propiedades específicas de cada clase (e.g., agua $\to$ neuronas 12--45, vegetación $\to$ neuronas 200--310).
  \item \textbf{Dimensionalidad efectiva}: El AFE (Sección 5.5) confirma que 6 factores capturan el 85.6\% de la varianza, indicando que el espacio de 512 dimensiones tiene una dimensionalidad intrínseca mucho menor. El SAE logra dispersión sin sacrificar la riqueza representacional: comprime la información en 6 dimensiones latentes mientras mantiene la capacidad de reconstrucción casi perfecta (MSE $= 10^{-6}$).
\end{enumerate}

\subsection{Conclusiones Parciales --- Situación 2}

\begin{enumerate}[nosep, leftmargin=*, itemsep=2pt]
  \item \textbf{Modelo funcional con KPIs cumplidos}: El modelo GeoVision-CLIP + SAE alcanza Recall@1 de 0.4765 (KPI $\geq$ 0.45), Recall@5 de 1.00 ($\geq$ 0.70), sparsity de 0.94 ($\geq$ 0.70) y MSE de reconstrucción de $10^{-6}$ ($\leq$ 0.05), superando los cuatro umbrales mínimos.
  \item \textbf{SAE como regularizador efectivo}: La combinación de dispersión L1 y normalización del decoder produce embeddings con alta separabilidad inter-clase y dimensionalidad efectiva baja (6 factores explican 85.6\% de varianza). La ablación confirma que el SAE aporta 2.44 puntos porcentuales al Recall@1.
  \item \textbf{LoRA previene sobreajuste}: Con solo 2,263 tiles, el fine-tuning completo colapsa el modelo (Recall@1 cae a 0.4392). LoRA (rank 8, 1.7M params) actúa como regularizador implícito limitando la capacidad de memorización.
  \item \textbf{Estructura factorial validada}: El AFC confirma un modelo de 6 factores con índices excelentes (RMSEA=0.042, CFI=0.999), proporcionando evidencia psicométrica rigurosa de que los embeddings SAE poseen una organización latente bien definida.
  \item \textbf{Base sólida para Situación 3}: Los embeddings dispersos y estructurados del SAE constituyen la entrada para el modelo Conv3D de predicción espacio-temporal, donde la baja dimensionalidad efectiva (6 factores) favorece la generalización con pocos datos de entrenamiento (1,403 secuencias).
\end{enumerate}

\vspace{0.1cm}
\noindent\rule{\columnwidth}{0.3pt}
\vspace{0.15cm}

% ════════════════════════════════════════════════════════════
\begin{thebibliography}{99}

\bibitem{veefkind}
\scriptsize J.P. Veefkind et al., ``TROPOMI on the ESA Sentinel-5 Precursor: A GMES mission for global observations of the atmospheric composition,'' \emph{Remote Sens. Environ.}, vol. 120, pp. 70--83, 2012.

\bibitem{vangeffen}
\scriptsize J. van Geffen et al., ``S5P TROPOMI NO$_2$ retrieval: impact of version v2.2 improvements,'' \emph{Atmos. Meas. Tech.}, vol. 15, pp. 2037--2060, 2022.

\bibitem{theys}
\scriptsize N. Theys et al., ``Sulfur dioxide retrievals from TROPOMI onboard Sentinel-5 Precursor,'' \emph{Atmos. Meas. Tech.}, vol. 10, pp. 119--153, 2017.

\bibitem{remoteclip}
\scriptsize F. Liu et al., ``RemoteCLIP: A Vision Language Foundation Model for Remote Sensing,'' \emph{IEEE Trans. Geosci. Remote Sens.}, vol. 62, pp. 1--16, 2024.

\bibitem{cressie}
\scriptsize N. Cressie \& C.K. Wikle, \emph{Statistics for Spatio-Temporal Data}. Hoboken, NJ: Wiley, 2011.

\bibitem{anselin}
\scriptsize L. Anselin, ``Local Indicators of Spatial Association---LISA,'' \emph{Geogr. Anal.}, vol. 27, no. 2, pp. 93--115, 1995.

\bibitem{templeton}
\scriptsize T. Templeton et al., ``Sparse Autoencoders for Mechanistic Interpretability,'' Anthropic Technical Report, 2023.

\end{thebibliography}

\label{LastPage}

\end{document}

